\scanpage{173}
\margin{$\circ\!\!-\!\!-\quad \omega_\ell(q)={}$ordre de $q\bmod\ell$ dans $(\mathbb{Z}/\ell\mathbb{Z})^{*}$}

\diagram{\begin{tikzcd}[ampersand replacement=\&,column sep=large]
k=\mathbb{F}_q \arrow[r,dash,"\pi\simeq\widehat{\mathbb{Z}}"'] \& \bar{k}
\end{tikzcd}}
engendré par $f=\mathrm{frob}^{q}$\qquad $M$ un $\mathbb{Z}_\ell$-module de t.f. sur le quel $\pi$ opère continûment.
\[
\begin{aligned}
H^0(\pi,M)&=M^\pi=M^{(1-f)}=\mathrm{Ker}\,(1-f)_M\\
H^1(\pi,M)&=M_{\deleted{\unsure{f|i}}\pi}=M_{(1-f)}=\mathrm{Coker}\,(1-f)_M=M/(1-f)M\\
H^i(\pi,M)&=0\qquad\text{si }\ i\geq 2
\end{aligned}
\]

\[
H^0(\pi,\mathbb{Z}_\ell(n))=0\quad\text{si }\ n\ne0,\quad =\mathbb{Z}_\ell\quad\text{si }\ n=0
\]
\[
\begin{aligned}
H^1(\pi,\mathbb{Z}_\ell(n))&\underset{\text{\inserted{non canoniquement}}}{=}\mathbb{Z}_\ell/(1-q^n)\mathbb{Z}_\ell=\begin{cases}
\mathbb{Z}_\ell&\text{si }\ n=0\\
0&\text{si }\ \ell\nmid(1-q^n)\quad\text{i.e. }\deleted{\unsure{n\ne0}}\ \inserted{n\not\equiv0\ (\omega_\ell(q))}\\
\mathbb{Z}/\ell^{\unsure{v|\nu}_\ell(1-q^n)}&\text{si }\ \underbrace{\ell\mid1-q^n}_{\inserted{\text{i.e. }n\equiv0\ (\omega_\ell(q))}},\ n\ne0
\end{cases}
\end{aligned}
\]
\margin{$\circ\!\!-\!\!-$}
\[
\left\{\begin{aligned}
H^{i}(\pi,\mathbb{Q}_\ell(n))&=0\quad\text{si }\ n\ne0,\quad\text{quel que soit }i\\
H^0(\pi,\mathbb{Q}_\ell)&=\mathbb{Q}_\ell\deleted{\unsure{f|1}},\ H^1(\pi,\mathbb{Q}_\ell)=\mathbb{Q}_\ell
\end{aligned}\right.
\]

\[
\deleted{H^i(\pi,\underline{\mathbb{Z}}(n))}
\]
\[
\left\{\begin{aligned}
\mathbb{H}^i(\mathbb{F}_q,\underset{\sim}{\mathbb{Q}}(n))&=0\quad\text{si }\ n\ne0,\quad\text{quel que soit }i\\
\mathbb{H}^0(\mathbb{F}_q,\underset{\sim}{\mathbb{Q}})&=\mathbb{Q},\ \mathbb{H}^1(\mathbb{F}_q,\underset{\sim}{\mathbb{Q}})=\underset{\sim}{\mathbb{Q}}.
\end{aligned}\right.
\]
\[
\left\{\begin{aligned}
H^0(\mathbb{F}_q,\underset{\sim}{\deleted{\mathbb{Q}}\mathbb{Z}})&=\mathbb{Z}\unsure{,|{}_\ell}\ (\text{mod }p\text{-gpes }?),\quad H^1(\mathbb{F}_q,\underset{\sim}{\mathbb{Z}})=\mathbb{Z}\quad(\text{mod }p\text{-gpes }?)\\
H^0(\mathbb{F}_q,\underset{\sim}{\mathbb{Z}}(\underset{\substack{\vert\\\inserted{n\ne0}}}{n}))&=0\quad(\text{id}),\quad H^1(\mathbb{F}_q,\underset{\sim}{\mathbb{Z}}(\underset{\substack{\vert\\\inserted{n\ne0}}}{n}))=\mathbb{Z}/(1-q^n)\mathbb{Z}\quad(\text{mod }p\text{-gpes }?)
\end{aligned}\right.
\]

\noindent\rule{.76\linewidth}{.35pt}
\[
\begin{aligned}
\mathbb{H}^i(\mathrm{Spec}\,\mathbb{Z},\mathbb{Z}(1))&=\begin{cases}0&\text{\unsure{si|sf}}\deleted{\unsure{<}}\ i\ne3\\\mathbb{Z}&\text{si }\ i=3\end{cases}\\
\mathbb{H}^i(\mathrm{Spec}\,\mathbb{Z},\mathbb{Z})&=\begin{cases}0&\text{si }\ i\ne0\\\mathbb{Z}&\text{si }\ i=0\end{cases}
\end{aligned}
\]
\[
H^i(\mathrm{Spec}(\mathbb{Z}),\deleted{\unsure{\mathbb{Z}}}\mathbb{Q}(n))\xrightarrow{\sim} H^{\unsure{i|1}}(\deleted{\unsure{k}}U,\deleted{\unsure{\mathbb{Z}}}\mathbb{Q}(n))
\]
si $U$ ouvert de Zariski de $\mathrm{Spec}\,\mathbb{Z}$, et $n\ne0$ (tout $i$).

% archive sheet number: 172
